\scanpage{79}
% lines 1–6, split product
Si $P_i$ sous $L_i$, avec $\pi_i$, est scindé par $x_i$ $(i=1,2)$, alors $P_1*P_2$ est scindé par $x_1*x_2$. Sa $\unsure{\chi|\gamma}$-trivialisation est donnée par un élément $b$ de $\deleted{\unsure{M}}L_1\otimes L_2$, et on aura
\[
b=b_1\otimes b_2
\]
Il n'y a pas plus simple !
% lines 7–9, introduction to structural data
Ceci permet d'expliciter les données d'associativité, de commutativité, d'unité pour l'opération $*$. Mais de façon intrinsèque :
% lines 10–14, commutativity
1) \uline{Commutativité} L'isom \deleted{\unsure{d}}can. de commutativité
\[
\deleted{\unsure{P}}E_1*E_2\xrightarrow{\sim}E_2*E_1
\]
transforme $x_1*x_2$ en $x_2*x_1$, et $u_1\otimes u_2$ en $u_2\otimes u_1$, ce qui suffit à le caractériser.
% lines 15–20, associativity and beginning of unit
2) \uline{Associativité} L'iso d'ass.
\[
(E_1*E_2)*\deleted{\unsure{(}}E_3\xrightarrow{\sim}E_1*(E_2*E_3)
\]
transforme
\[
(x_1*x_2)*x_3\quad\text{en}\quad x_1*(x_2*x_3)
\]
et
\[
(\deleted{\unsure{x}}u_1\otimes u_2)\otimes u_3\quad\text{en}\quad u_1\otimes(u_{\unsure{2|1}}\otimes u_3),
\]
ce qui suffit à le caractériser.

3) \uline{Unité} : On \deleted{appelle} dit qu'une
% archive sheet number: 78
